%% LyX 2.4.4 created this file.  For more info, see https://www.lyx.org/.
%% Do not edit unless you really know what you are doing.
\documentclass[english]{article}
\usepackage[T1]{fontenc}
\usepackage[latin9]{inputenc}
\usepackage{enumitem}
\usepackage{amsmath}
\usepackage{amsthm}
\usepackage{amssymb}
\usepackage{geometry}
\geometry{verbose,tmargin=1in,bmargin=1in,lmargin=1in,rmargin=1in}

\makeatletter
%%%%%%%%%%%%%%%%%%%%%%%%%%%%%% Textclass specific LaTeX commands.
\numberwithin{equation}{section}
\numberwithin{figure}{section}
\newlength{\lyxlabelwidth}      % auxiliary length 

%%%%%%%%%%%%%%%%%%%%%%%%%%%%%% User specified LaTeX commands.
\usepackage{fancyhdr}
\usepackage{lastpage}
\pagestyle{fancy}
\usepackage{enumitem}
\usepackage{ifthen}
\everymath=\expandafter{\the\everymath\displaystyle}
%\DeclareMathSizes{10}{10}{10}{10}
\usepackage{multicol}

\usepackage{tikz}
\usetikzlibrary{patterns}
\usetikzlibrary{plotmarks}
\usepackage{pgfplots}
\pgfplotsset{compat=newest}

\usepackage{calc}
\usepackage[nomessages]{fp}

\usepackage{datetime}
% Added by lyx2lyx
\setlength{\parskip}{\smallskipamount}
\setlength{\parindent}{0pt}

\makeatother

\usepackage{babel}
\begin{document}
\renewcommand{\author}{Dr.\,James Ashe}

\newcommand{\classNum}{439}

\newcommand{\classSec}{A}

\newcommand{\nameType}{Name: \hspace{-4cm}}

\newcommand{\examType}{Homework 1}

\newcommand{\Date}{\formatdate{14}{1}{2014}} %%\AdvanceDate[12] \today

%\newdate{my_date}{\day}{\month}{\year}%   
%\AdvanceDate[-5] 
%\displaydate{my_date}

%%First page's header%%%%%%%%%%%%%%%%%%%%%%
\fancypagestyle{firststyle} %%header settings for first pagr
{
	\fancyhead[L]{MTH \classNum{}(\classSec{}) \author{}\\
	\textbf{\examType{}} \Date{}}
	\fancyhead[C]{\textbf{\nameType}}
	\setlength{\headsep}{14pt}
}
%%%%%%%%%%%%%%%%%%%%%%%%%%%%%%%%%%%%%%%%%%

%%Next page's header%%%%%%%%%%%%%%%%%%%%%%%
%\setlist{nolistsep} %eliminates space between list items
\lhead{MTH \classNum{}(\classSec{})} 
\chead{\textbf{\nameType}} 
\cfoot{\thepage{}~of~\pageref{LastPage}} 
\setlength{\headsep}{5pt} 
%%%%%%%%%%%%%%%%%%%%%%%%%%%%%%%%%%%%%%%%%%%

%%Various settings; see the preamble for other settings%%%%%
%%\setenumerate[0]{leftmargin=12pt,itemindent=0pt,labelwidth=0pt}
\renewcommand{\labelenumi}{\textbf{\arabic{enumi}.}}
\renewcommand{\labelenumii}{\textbf{\alph{enumii}.}}
\thispagestyle{firststyle}
%%%%%%%%%%%%%%%%%%%%%%%%%%%%%%%%%%%%%%%%%%

\textbf{Homework}
\begin{enumerate}
\item Show that $\mathbb{Z}_{6}$ is a ring with unity under addition and
multiplication. 
\item Show that the set $\sqrt{2}\mathbb{Z}=\{x\colon x=\sqrt{2}n\mbox{ for some \ensuremath{n\in\mathbb{Z}}}\}$
is a ring. 
\item Let $F$ be the set of all functions from $\mathbb{R}$ to $\mathbb{R}$.
Define addition and multiplication as 
\[
(f+g)(x)=f(x)+g(x)\mbox{ and }(fg)(x)=f(x)g(x).
\]
Show that $F$ is a commutative ring with identity. 
\item Prove that the ring above, $F$, is not an integral domain.

\begin{itemize}
\item [Hint]Here are two \emph{indicator functions} 
\[
\mathbf{1}_{[0,\infty)}(x)=\begin{cases}
1 & \mbox{if }x\in[0,\infty)\\
0 & \mbox{if }x\notin[0,\infty)
\end{cases}\mbox{ and }\mathbf{1}_{(-\infty,0]}(x)=\begin{cases}
1 & \mbox{if }x\in(-\infty,0]\\
0 & \mbox{if }x\notin(-\infty,0]
\end{cases}
\]
\\
 The indicator function is a function that has an output of either
$1$ or $0$. For example, $\mathbf{1}_{[0,\infty)}(2)=1$ and $\mathbf{1}_{[0,\infty)}(-2)=1$.
Now consider the functions $f(x)=x\mathbf{1}_{[0,\infty)}$ and $g(x)=x\mathbf{1}_{(-\infty,0]}$.
\end{itemize}
\item Let $R$ and $S$ be rings. Define addition and multiplication on
the Cartesian product $R\times S$ by, 
\[
(r,s)+(r',s')=(r+r',s+s')\mbox{ and }(r,s)(r',s')=(rr',ss')
\]
for $r,r'\in R$ and $s,s'\in S$. Prove that $R\times S$ is a ring. 
\item Show that the set $k\mathbb{Z}=\{x\colon x=kn\mbox{ for some \ensuremath{n\in\mathbb{Z}}}\}$
for any integer $k$ is a subring of $\mathbb{Z}$. 
\item Prove that $\mathbb{Q}$ is an integral domain. 
\item Prove that $\mathbb{Z}_{6}$ is not an integral domain. 
\item Let $a$ and $b$ be two nonzero elements in a ring $R$. Prove that
in general, $ax=b$ does not have a unique solution.

\begin{itemize}
\item [Hint]It should be understood that $x$ must be in $R$. An example
is all that is needed. You can find an example that has two solutions,
or find an example with no solution. 
\end{itemize}
\item Prove that $\mathbb{Z}_{n}$ is not an integral domain (under the
usual addition and multiplication) if $n$ is not a prime ($n$ is
assumed to be a positive integer). 

\begin{itemize}
\item [Hint]If $n$ is not prime then $n=mq$ for two integers $1<m,q<n$.
Looking at the example seen in exercise might help. 
\end{itemize}
\end{enumerate}
\textbf{Presentation Problems} 
\begin{enumerate}
\item Prove that the intersection of any two subrings of a ring $R$ is
also a subring of $R$. 
\item Give an example of a ring with elements $a$ and $b$ with properties
that $ab=0$ but $ba\neq0$. 
\item Let $m$ and $n$ be positive integers with least common multiple
$k$. Proves that $m\mathbb{Z}\cap n\mathbb{Z}=k\mathbb{Z}$. 
\item Prove that a ring that is cyclic under addition is commutative. 
\item Prove that if a ring element has a multiplicative inverse, it is unique. 
\item Prove that a field is an integral domain. 
\item Prove that $\mathbb{Z}_{p}$ is an integral domain if $p$ is prime
(you can assume that it is a commutative ring with identity). 
\item Prove that $\mathbb{Z}_{p}$ is a field.
\end{enumerate}

\end{document}
